%!TEX root=thesis.tex
Pengimejan Resonans Magnetik Jantung (CMR) digunakan secara meluas untuk penilaian kuantitatif anatomi dan fungsi jantung. Dalam analisis jantung automatik, segmentasi yang tepat bagi ventrikel kanan (RV), miokardium (MYO), dan ventrikel kiri (LV) merupakan langkah penting sebelum ukuran fungsi seperti isipadu akhir diastol, isipadu akhir sistol, jisim miokardium, dan pecahan ejeksi boleh diperoleh. Walaupun seni bina berasaskan U-Net masih menjadi asas yang kukuh dalam segmentasi imej perubatan, pelaksanaan praktikal juga perlu mempertimbangkan saiz model, kos inferens, dan protokol pengesahan yang boleh diulang. Tesis ini mengkaji Nano-Mamba U-Net, iaitu seni bina 3D U-Net ringan yang diinspirasikan oleh Mamba untuk segmentasi MRI jantung. Model ini mengekalkan struktur pengekod-penyahkod dan sambungan langkau U-Net, sambil memasukkan lapisan bottleneck pemprosesan jujukan yang padat pada perwakilan laten termampat. Lapisan ini menyusun ciri 3D kepada jujukan satu dimensi, menggunakan modul jujukan berpagar ringan yang diinspirasikan oleh pemodelan ruang keadaan gaya Mamba, dan membentuk semula output kepada bentuk volumetrik. Oleh itu, pelaksanaan ini diterangkan sebagai Mamba-inspired dan bukannya pelaksanaan penuh algoritma selective scan asal.

Model yang dicadangkan dinilai menggunakan set data Automated Cardiac Diagnosis Challenge (ACDC) dengan protokol pengesahan berasaskan pesakit. Eksperimen akhir menggunakan 80 pesakit untuk latihan dan 20 pesakit untuk pengesahan, bersamaan dengan 160 kes latihan dan 40 kes pengesahan. Semua model perbandingan dilatih menggunakan prapemprosesan dan prosedur pengesahan yang sama, dan checkpoint terbaik bagi setiap model dipilih berdasarkan min Dice pengesahan. Nano-Mamba U-Net mencapai min Dice Similarity Coefficient (DSC) pengesahan sebanyak 84.78\% dengan 1.46 juta parameter boleh latih. Model ini mengatasi garis dasar 3D U-Net, yang mencapai 80.83\% min DSC dengan 4.81 juta parameter. Walau bagaimanapun, SegResNet16 mencapai min DSC tertinggi sebanyak 86.70\%, manakala ablation No-Mamba mencapai 85.64\%. Keputusan ini menunjukkan bahawa Nano-Mamba U-Net tidak harus didakwa sebagai model paling tepat berdasarkan Dice. Sebaliknya, sumbangan utamanya ialah keseimbangan antara ketepatan dan kecekapan: model ini memberikan prestasi pengesahan yang kompetitif dengan bilangan parameter paling kecil dalam kalangan model yang dinilai. Tesis ini menyimpulkan bahawa bottleneck ringan yang diinspirasikan oleh pemodelan jujukan berpotensi untuk segmentasi MRI jantung, namun kajian lanjut masih diperlukan untuk penilaian set ujian rasmi, analisis khusus patologi, pengesahan luaran, dan pemodelan gerakan jantung 4D sepenuhnya.

\vspace{1em}
\noindent
\textbf{Kata Kunci:} Segmentasi MRI Jantung, Pembelajaran Mendalam, Seni Bina Diinspirasikan Mamba, Model Ruang Keadaan, Keseimbangan Ketepatan-Kecekapan.
